\section{Background and Methodology}
\label{sec:methodology}

Argo floats park and drift at 1000m depth between profiling cycles.
Their locations are determined by currents at that depth, not by experimental design.
If 1000m currents are correlated with upper-ocean temperature or salinity, then Argo-based estimates of quantities like space-time averaged ocean heat content (OHC) could be biased.
This is widely suspected to be negligible, but the magnitude of this bias has not yet been bounded statistically.

This section sets up the estimand, the sampling mechanism, and the resulting sampling-induced bias.
Three logically distinct sources of randomness appear, each on its own probability space: a \emph{field} space carrying the ocean heat content (OHC) field, which---following model-based geostatistics and the hierarchical Bayesian OHC model of \citet{baugh2022} on which this work builds---we treat as a \emph{stochastic process}; a \emph{sampling} space carrying the Argo float point process; and a \emph{population} space on which a random space--time coordinate is drawn and against which the target average is defined.
Our analysis of the sampling bias is carried out \emph{conditional} on a realization of the field and on the sampling design; conditional on these, the field and the sampling weight are fixed functions and the only operative randomness is the target coordinate.
This conditioning is what allows a distribution-free bound to coexist with, and indeed validate, a stochastic-field model rather than contradict it.

\subsection{The space--time domain and the target measure}
\label{sec:domain}

Fix a spatial region of interest $A \subset S^{2}$, taken to be a
Borel subset of the sphere with finite, positive surface area, and a
study period $\T \subset \R$, a bounded interval.

\begin{definition}[Space--time domain]
\label{def:domain}
The \emph{space--time domain} is the product set
$\X := A \times \T$, equipped with the product $\sigma$-algebra
$\F := \B(A)\otimes\B(\T)$, where $\B(A)$ and $\B(\T)$ are the Borel
$\sigma$-algebras on $A$ and $\T$. A generic element of $\X$ is written
$x = (s,t)$, with spatial location $s \in A$ and time $t \in \T$.
\end{definition}

To speak of a space--time \emph{average} we specify the weighting by
which space and time are aggregated; this is the role of the reference
measure $\nu$, which encodes the target notion of ``uniform over the
domain.''

\begin{definition}[Reference space--time measure]
\label{def:nu}
Let $|\cdot|$ denote surface (area) measure on $(A,\B(A))$, and define
the \emph{normalized area measure} $\alpha(E) := |E|/|A|$,
$E \in \B(A)$, a probability measure on $(A,\B(A))$. Let $\tau$ be a
fixed probability measure on $(\T,\B(\T))$; unless stated otherwise
$\tau$ is normalized Lebesgue measure, so time is weighted uniformly.
The \emph{reference space--time measure} is the product probability
measure
\[
  \nu := \alpha \otimes \tau \qquad\text{on}\qquad (\X,\F).
\]
\end{definition}

Because $\alpha$ and $\tau$ are probability measures, so is $\nu$;
hence $(\X,\F,\nu)$ is a probability space, the \emph{population} (or
\emph{target}) space. It carries no information about how the ocean is
sampled or about the value of the field; it encodes only the averaging
weights that define the quantity we wish to estimate.

\subsection{The target-population random point}
\label{sec:random-point}

The estimand is a deterministic integral against $\nu$. To make the
covariance and inner-product machinery of Section~\ref{sec:bias-decomp}
available, we re-express that integral as an expectation by introducing
a random space--time coordinate whose distribution is exactly $\nu$. This
coordinate lives on the population space and is unrelated to both the
field space of Section~\ref{sec:field} and the sampling space of
Section~\ref{sec:sampling}.

\begin{definition}[Target-population random point]
\label{def:X}
Let $X = (S,T)$ be the \emph{coordinate map} on $(\X,\F,\nu)$, i.e. the
identity $X : \X \to \X$, $X(s,t) = (s,t)$. Its distribution is the reference
measure by construction, $X_{\#}\nu = \nu$; equivalently
$\Prob(X \in E) = \nu(E)$ for every $E \in \F$. We write $X \sim \nu$.
\end{definition}

Because $\nu = \alpha\otimes\tau$ is a product, $S \sim \alpha$
(uniform by surface area on $A$) and $T \sim \tau$ are independent.
Concretely, $X$ is a point obtained by sampling a location uniformly by
area and, independently, a time from the study period. The single
identity used repeatedly is the change-of-variables formula
\begin{equation}
  \E_\nu\big[g(X)\big] = \int_\X g \dd\nu
  \qquad\text{for every $\F$-measurable } g \in L^1(\nu).
  \label{eq:cov}
\end{equation}
$X \sim \nu$ is a representation device, not a physical model: it lets
the target average be read as $\E_\nu[\,\cdot\,]$. It is not a float
location and is unrelated to the sampling process.

\subsection{The OHC field as a random field}
\label{sec:field}

Model-based geostatistics represents the unknown spatial phenomenon as
a stochastic process rather than a fixed surface. We adopt that view so
that the present analysis is consistent with the hierarchical Bayesian
OHC model of \citet{baugh2022} on which it builds: the OHC field is
random, carried on its own probability space.

\begin{definition}[OHC random field]
\label{def:Y}
Let $(\Omega_Y,\mathcal{A}_Y,\Prob_Y)$ be a probability space, the
\emph{field space}. The \emph{OHC field} is a random element
\[
  Y : \Omega_Y \to L^2(\nu), \qquad \omega_Y \mapsto Y(\omega_Y,\cdot),
\]
jointly measurable in $(\omega_Y,x)$, with $Y(\omega_Y,\cdot)\in L^2(\nu)$
for $\Prob_Y$-almost every $\omega_Y$ (and, for the OHC application,
with continuous sample paths, so that point evaluations $Y(x_i)$ at
float locations are well defined). Its distribution $\Prob_Y\circ Y^{-1}$ is the
\emph{field distribution}; in the OHC application this distribution is modeled
as a (nonstationary) Gaussian process, and the posterior of \citet{baugh2022} given the Argo profiles approximates its conditional
distribution given the data.
\end{definition}

The bias we study is a property of the \emph{realized} ocean, so we
condition on a realization of the field throughout.

\begin{definition}[Conditioning on the realized field]
\label{def:realization}
Fix $\omega_Y\in\Omega_Y$ and write $y := Y(\omega_Y,\cdot)\in L^2(\nu)$
for the corresponding realization. Because the study concerns a single
physical ocean over the period $\T$, we identify $y$ with the true
realized OHC field and treat an independent reanalysis (GLORYS12) as a
proxy for $y$. All quantities in
Sections~\ref{sec:conditioning}--\ref{sec:bias-decomp} are defined
conditional on the event $\{Y = y\}$; conditional on it, $y$ is a fixed
element of $L^2(\nu)$.
\end{definition}

Composing $y$ with the coordinate map yields the real-valued random
variable $y(X)$ on the population space, whose expectation is the
quantity of interest.

\begin{definition}[Estimand: realized regional-mean OHC]
\label{def:mu}
The \emph{realized target mean OHC} is
\[
  \mu(y) := \int_\X y \dd\nu = \E_\nu\big[y(X)\big].
\]
Viewed as a functional of the random field, $\mu(Y)$ is itself a random
variable on the field space, whose distribution the OHC model
characterizes and for which \citet{baugh2022} report credible
intervals; conditional on $\{Y=y\}$ it is a fixed scalar. Write
\[
  \sigma_Y(y) := \|\,y - \mu(y)\,\|_{L^2(\nu)}
\]
for the space--time standard deviation of the realized field, finite
since $y\in L^2(\nu)$.
\end{definition}

\begin{remark}[Estimating $\sigma_Y$ without circularity]
\label{rem:glorys}
Estimating $\sigma_Y(y)$ requires the realized field $y$ over all of
$\X$, including unsampled space--time. Using the Argo-fit OHC posterior
for this purpose would estimate the sampling bias of a model from the
same preferentially sampled data whose ignorability is in question. We
therefore take an independent reanalysis (GLORYS12), which assimilates
information beyond Argo, as the primary proxy for $y$, and report any
posterior-based estimate only as a secondary sensitivity check.
\end{remark}

\subsection{The sampling process and the sampling measure}
\label{sec:sampling}

The observations are produced by Argo floats, realizing a
spatio-temporal point process on a probability space separate from both
the population space and the field space.

\begin{definition}[Observation point process]
\label{def:N}
Let $(\Omega,\mathcal{A},\Prob)$ be a probability space, the
\emph{sampling space}. The float profiles are a point process $N$ on
$(\X,\F)$ over this space: a measurable map $N:\Omega\to\mathsf{M}(\X)$,
$\omega\mapsto N(\omega,\cdot)$, where $\mathsf{M}(\X)$ is the space of
counting measures on $(\X,\F)$. For $E\in\F$, $N(\omega,E)$ is the
number of profiles in $E$ for realization $\omega$.
\end{definition}

\begin{definition}[Intensity measure]
\label{def:Lambda}
The \emph{intensity measure} of $N$ is the expected counting measure
\[
  \Lambda(E) := \E_N\big[N(E)\big] = \int_\Omega N(\omega,E)\dd\Prob(\omega),
  \qquad E \in \F .
\]
\end{definition}

$\Lambda$ is a non-random measure on the population space: the
expectation over $\Omega$ has integrated out the realization.

\begin{assumption}[Finite, dominated intensity]
\label{ass:intensity}
The total intensity is finite and positive, $0<\Lambda(\X)<\infty$,
and $\Lambda \ll \nu$.
\end{assumption}

\begin{definition}[Sampling probability measure]
\label{def:pi}
The \emph{sampling probability measure} is $\pi := \Lambda/\Lambda(\X)$
on $(\X,\F)$; $\pi(E)$ is the probability that a profile drawn at random
from the sampling process lies in $E$. When a sampled location is
treated as random we write $X' \sim \pi$.
\end{definition}

\subsection{The sampling weight}
\label{sec:weight}

The comparison between how the ocean is sampled ($\pi$) and how the
target average weights it ($\nu$) is a single density.

\begin{definition}[Sampling weight]
\label{def:w}
Under Assumption~\ref{ass:intensity}, $\pi \ll \nu$, and the
\emph{sampling weight} is the Radon--Nikodym derivative
$w := \dd\pi/\dd\nu$: the $\nu$-almost-everywhere unique, non-negative,
$\F$-measurable function with $\pi(E)=\int_E w\dd\nu$ for all $E\in\F$.
As an element of $L^1(\nu)$ it is mean-one,
\begin{equation}
  \E_\nu[w] = \int_\X w \dd\nu = \pi(\X) = 1,
  \label{eq:mean-one}
\end{equation}
so $w-1$ is centered. For the second-moment bound we further assume
$w\in L^2(\nu)$, and write $\sigma_w := \|w-1\|_{L^2(\nu)}$.
\end{definition}

Two structural points fix the status of $w$. First, $w$ is a function
on the population space, built from $\Lambda=\E_N[N]$, in which the
float randomness $\omega\in\Omega$ has already been averaged out; hence
$w$ does not depend on $\omega$ and is deterministic at the level of the
sampling design. Second, $w$ is an importance ratio: for every
$g\in L^1(\pi)$,
\begin{equation}
  \E_\pi\big[g(X')\big] = \int_\X g\,w\dd\nu = \E_\nu\big[w(X)\,g(X)\big],
  \label{eq:importance}
\end{equation}
so $w(X)$ reweights a target draw $X\sim\nu$ into a sampled draw
$X'\sim\pi$. Non-preferential sampling is $w\equiv 1$ ($\pi=\nu$), for
which $\sigma_w=0$. Because $\E_\nu[w]=1$, the scalar $\sigma_w$ is the
standard deviation, equivalently the coefficient of variation, of the
mean-one variable $w(X)$: a summary of how inhomogeneous the sampling
is.

\begin{remark}[Estimated weights]
\label{rem:hat-w}
In practice $\pi$, hence $w$, are unknown and estimated from the
realized float locations (kernel density estimation, a log-Gaussian Cox
process, or a Gaussian-process intensity model). The estimate $\hat w$
is a random element of $L^2(\nu)$, measurable with respect to
$\sigma(N)$ on the sampling space. Conditional on $\sigma(N)$---i.e.
given the observed locations---$\hat w$ is fixed, and every bound below
is a deterministic functional of $\hat w$. Sampling variability of
$\hat w$ about $w$ is a separate estimation error, propagated
separately.
\end{remark}

\subsection{Sources of randomness and the conditioning stance}
\label{sec:conditioning}

Three probability spaces have now been introduced, each carrying a
distinct kind of randomness:
\begin{itemize}
  \item the \emph{field space} $(\Omega_Y,\mathcal{A}_Y,\Prob_Y)$,
        carrying the random OHC field $Y$ (Definition~\ref{def:Y});
  \item the \emph{sampling space} $(\Omega,\mathcal{A},\Prob)$,
        carrying the float point process $N$, hence any estimate
        $\hat w$ (Definitions~\ref{def:N} and~\ref{rem:hat-w});
  \item the \emph{population space} $(\X,\F,\nu)$, carrying the target
        coordinate $X$ (Definition~\ref{def:X}).
\end{itemize}

The bias analysis is conditional. We condition on a realization
$\{Y=y\}$ of the field and on the sampling design (equivalently, on the
intensity $\Lambda$, hence on the population weight $w$). Conditional on
these, $y$ and $w$ are fixed elements of $L^2(\nu)$, and the only
operative randomness in every identity and bound of
Section~\ref{sec:bias-decomp} is the target coordinate $X\sim\nu$; all
covariance and inner-product operations there are taken over $X$ alone.

\begin{remark}[No joint model of field and locations is required]
\label{rem:no-joint}
Conditioning on $\{Y=y\}$ sidesteps any need to specify the joint distribution
of the field and the float locations---the object a generative
preferential-sampling model (e.g. that of \citet{diggle2010}) would
posit through an intensity that depends on the field. Our bound holds
for the realized ocean and the realized design without assuming a
mechanism linking them, which is the source of its distribution-free
character.
\end{remark}

\begin{remark}[Coherence with the conditional OHC model]
\label{rem:coherence}
Treating $Y$ as stochastic and then conditioning on a realization is
the standard device of model-based geostatistics, and it is precisely
the stance of the hierarchical OHC model, which conditions on the float
locations. The validity of that model rests on sampling being
ignorable for the field. The present analysis quantifies the bias that
non-ignorability could induce; when the bound and the estimated bias
are small, it \emph{validates} the conditional model rather than
contradicting it.
\end{remark}

\subsection{The design mean and the bias decomposition}
\label{sec:bias-decomp}

Throughout this subsection all quantities are conditional on the
realized field $\{Y=y\}$ and on the design, as set out in
Section~\ref{sec:conditioning}.

\begin{definition}[Design mean]
\label{def:mu-w}
The \emph{design mean} is the realized field averaged under the
sampling measure,
\[
  \mu_w(y) := \int_\X y \dd\pi = \E_\pi\big[y(X')\big]
  = \int_\X w\,y \dd\nu = \E_\nu\big[w(X)\,y(X)\big],
\]
using the importance identity~\eqref{eq:importance}.
\end{definition}

Let $\hat\mu$ denote any estimator of $\mu(y)$ built from the observed
profiles (for example a kriging/Gaussian-process interpolant). Its bias
decomposes into a sampling term and a reconstruction term.

\begin{definition}[Bias decomposition]
\label{def:bias}
The bias of $\hat\mu$ decomposes as
\[
  B(y) := \hat\mu - \mu(y)
  = \underbrace{\big(\mu_w(y) - \mu(y)\big)}_{\textstyle B_{\mathrm{samp}}(y)}
  + \underbrace{\big(\hat\mu - \mu_w(y)\big)}_{\textstyle B_{\mathrm{pred}}(y)},
\]
where $B_{\mathrm{samp}}(y)$ is the \emph{preferential-sampling bias}
and $B_{\mathrm{pred}}(y)$ is the \emph{prediction bias}. The split is a
modeling convention: $B_{\mathrm{samp}}$ is defined estimator-free, so
any correction an interpolator makes for uneven coverage is charged to
$B_{\mathrm{pred}}$. Because $\pi$ is built from the intensity
$\Lambda=\E_N[N]$, $B_{\mathrm{samp}}$ is the systematic
(intensity-level) sampling bias; realization-to-realization fluctuation
of $N$ contributes to $B_{\mathrm{pred}}$, not to $B_{\mathrm{samp}}$.
\end{definition}

The sampling term admits a covariance representation on the population
space, the starting point for the bound.

\begin{proposition}[Sampling-bias covariance identity]
\label{prop:cov-identity}
Conditional on $\{Y=y\}$ and the design, and under
Assumption~\ref{ass:intensity} with $y,w\in L^2(\nu)$,
\[
  B_{\mathrm{samp}}(y) = \mu_w(y) - \mu(y)
  = \E_\nu\big[(w-1)\,y\big]
  = \Cov_\nu\big(w(X),\,y(X)\big).
\]
\end{proposition}

\begin{proof}
By Definitions~\ref{def:mu} and~\ref{def:mu-w},
$\mu_w(y)-\mu(y)=\E_\nu[wy]-\E_\nu[y]=\E_\nu[(w-1)y]$. Expanding the
covariance about $\mu(y)$ and using~\eqref{eq:mean-one},
\[
  \Cov_\nu(w,y)
  = \E_\nu\big[(w-1)(y-\mu(y))\big]
  = \E_\nu\big[(w-1)y\big] - \mu(y)\,\E_\nu[w-1]
  = \E_\nu\big[(w-1)y\big],
\]
since $\E_\nu[w-1]=0$. Both moments are finite by Cauchy--Schwarz, as
$y,w-1\in L^2(\nu)$.
\end{proof}

\begin{remark}[Two distinct covariances---do not conflate]
\label{rem:two-cov}
The covariance in Proposition~\ref{prop:cov-identity} is taken over the
target coordinate $X\sim\nu$ with the field realization $y$ held
\emph{fixed}; it is a scalar---the space--time covariance between the
weight and the field across the domain---and it \emph{is} the
conditional sampling bias. It must not be conflated with the field's
covariance function
$k(x,x') = \Cov_{\Prob_Y}\!\big(Y(x),Y(x')\big)$, taken over field
randomness on $\Omega_Y$: the (nonstationary) covariance kernel of the
OHC model. The two objects live on different probability spaces and
play different roles. Because we condition on $y$, the kernel $k$ does
\emph{not} enter the $L^2$ bound; it would reappear only if one
measured field smoothness in the RKHS of $k$, or averaged the bias over
the field distribution (Remark~\ref{rem:expected}).
\end{remark}

\begin{remark}[No sampling bias under uncorrelated design]
\label{rem:zero-bias}
If, for the realized field, the weight and the field are uncorrelated
under $\nu$---in particular whenever $w\equiv1$---then
$\Cov_\nu(w,y)=0$ and $B_{\mathrm{samp}}(y)=0$. Preferential-sampling
bias arises precisely from co-variation between where the floats
concentrate ($w$) and the value of the field ($y$).
\end{remark}

The covariance representation yields the paper's central bound by a
single application of Cauchy--Schwarz in $L^2(\nu)$.

\begin{proposition}[Distribution-free $L^2$ bound]
\label{prop:l2-bound}
Under the conditions of Proposition~\ref{prop:cov-identity},
\[
  \big|B_{\mathrm{samp}}(y)\big|
  = \big|\Cov_\nu(w,y)\big|
  = |\rho_{w,y}|\,\sigma_w\,\sigma_Y(y)
  \;\le\; \sigma_w\,\sigma_Y(y),
\]
where $\rho_{w,y}\in[-1,1]$ is the $\nu$-correlation between $w(X)$ and
$y(X)$. The bound requires only finite second moments and imposes no
smoothness assumption on the field.
\end{proposition}

The worst case $|\rho_{w,y}|=1$ is attained only when $y-\mu(y)\propto
w-1$. The realized correlation $\rho_{w,y}$---estimable, with
$\sigma_w$ and $\sigma_Y(y)$, by quadrature against $\nu$ using the
GLORYS12 proxy for $y$ and an estimate $\hat w$ of the weight---is the
scalar that separates the pessimistic envelope $\sigma_w\sigma_Y(y)$
from the actual sampling bias, and is the empirical target of the paper.

\begin{remark}[Expected bias over the field distribution; kernel methods]
\label{rem:expected}
Averaging over the field distribution gives the expected sampling bias
$\E_{\Prob_Y}\!\big[B_{\mathrm{samp}}(Y)\big]
= \E_{\Prob_Y}\!\big[\Cov_\nu(w,Y)\big]$, which, unlike the conditional
identity, depends on the joint distribution of field and design and is the
natural object of a generative analysis. Likewise, replacing the
distribution-free class $L^2(\nu)$ by a smoothness ball in the RKHS of
$k$ yields a tighter, kernel-dependent bound. We pursue neither here:
the paper's aim is a distribution-free bound on the realized bias, and
an estimate of its magnitude, for which the conditional $L^2$
formulation of Proposition~\ref{prop:l2-bound} suffices. Both
extensions are deferred to future work, where the choice of kernel
would be tied to the OHC covariance model.
\end{remark}
